% Le mono jojoi
%=============================================
% 8.1. Explicar que también se utilizó búsqueda por grilla como apoyo visual
%=============================================

\begin{frame}{Algoritmo Numérico: Búsqueda por Grilla}

    \begin{columns}[T]
        
        %=======================================================
        % COLUMNA IZQUIERDA: BLOQUE VERDE (Pseudocódigo)
        %=======================================================
        \column{0.63\textwidth}
        \begin{exampleblock}{Lógica del Algoritmo de Exploración}
            \scriptsize % Letra pequeña para que las fórmulas largas quepan sin problema
            \textbf{Inicialización:} $E_{\text{min}} \leftarrow \infty$ \\[0.15cm]
            \begin{tabbing}
                \quad \= \quad \= \quad \= \quad \= \kill
                $\mathbf{para}$ cada $x$ dentro del rango sugerido $\mathbf{hacer}$ \\
                \> $\mathbf{para}$ cada $y$ dentro del rango sugerido $\mathbf{hacer}$ \\
                \> \> $\mathbf{para}$ cada $z$ dentro del rango sugerido $\mathbf{hacer}$ \\[0.1cm]
                \> \> \> 1. $R_i = \sqrt{(x_i - x)^2 + (y_i - y)^2 + (z_i - z)^2}$ \\
                \> \> \> 2. $A_{pred,i} = A_0 \frac{e^{-R_i}}{R_i}$ \\
                \> \> \> 3. $E(x,y,z) = \sum_{i=1}^{n} (A_{obs,i} - A_{pred,i})^2$ \\[0.1cm]
                \> \> \> $\mathbf{si}$ $E(x,y,z) < E_{\text{min}}$ $\mathbf{entonces}$ \\
                \> \> \> \> $E_{\text{min}} \leftarrow E(x,y,z)$ \\
                \> \> \> \> $\text{mejor\_punto} \leftarrow (x,y,z)$
            \end{tabbing}
        \end{exampleblock}

        %=======================================================
        % COLUMNA DERECHA: BLOQUES AZULES APILADOS
        %=======================================================
        \column{0.34\textwidth}
        
        \begin{block}{Procesamiento Espacial}
            \scriptsize
            \begin{itemize}
                \item \textbf{Discretización:} Evaluación en mallas tridimensionales.
                \vspace{0.1cm}
                \item \textbf{Métrica:} Validación simultánea en la red de sensores.
            \end{itemize}
        \end{block}

        \vspace{0.3cm}

        \begin{block}{Soporte Conceptual}
            \scriptsize
            \begin{itemize}
                \item \textbf{Topografía:} Mapea la superficie de error.
                \vspace{0.1cm}
                \item \textbf{Garantía:} Valida visualmente la existencia del mínimo global.
            \end{itemize}
        \end{block}

    \end{columns}

    \vspace{0.3cm}
    \begin{center}
        \footnotesize
        \textbf{La estructura iterativa asegura un rastreo exhaustivo del hipocentro.}
    \end{center}

\end{frame}